% ===========================================================================
% SOURCES: docs/controller_design.md (Greek - the roofline and the parallelism
%            study are here; reuse the Greek text for the extended summary
%            rather than translating twice)
%          docs/references/MobileNetOrigPaper.pdf, 1712.05877v1.pdf (Jacob),
%            wp486-deep-learning-int8.pdf, ug579-ultrascale-dsp.pdf,
%            ug1085-zynq-ultrascale-trm.pdf
%
% DISCIPLINE FOR THIS CHAPTER: background earns its place only where a later
% chapter depends on it. Every subsection below should end up cited from
% chapters 5-10. If one is not, cut it - a thesis is not a textbook.
% ===========================================================================
\chapter{Background and Technical Foundations}
\label{ch:background}

\section{Convolutional Neural Networks}
\label{sec:cnn}

\subsection{Convolutional Layers}
\label{sec:conv-layers}
% Keep short. The loop nest defined here is reused verbatim in
% \cref{sec:tiling} and \cref{ch:architecture}, so fix the notation now:
% OH, OW, IC, OC, KH, KW, stride, padding - and use it everywhere after.

\subsection{Training and Transfer Learning}
\label{sec:transfer}
% ImageNet pretraining [deng2009imagenet] -> fine-tuning on 4 grape classes.
% Only as much as \cref{sec:training} needs.

\section{The MobileNetV2 Architecture}
\label{sec:mobilenetv2}

\subsection{Depthwise Separable Convolution}
\label{sec:dwsep}
% [howard2017mobilenets]. Give the MAC-count argument, because it is what
% justifies having two separate engines in \cref{ch:architecture}: the
% pointwise convolutions carry \thPwMacPct\% of the MACs and the depthwise
% ones \thDwMacPct\%, which is why they get 512 MACs/cycle and 16 channels
% respectively rather than equal budgets.

\subsection{Inverted Residual Blocks and Linear Bottlenecks}
\label{sec:inverted-residual}
% [sandler2018mobilenetv2]. The residual connection is what forces the
% requantisation question in \cref{sec:residual-add}; flag that forward.

\subsection{Network Structure}
\label{sec:mbv2-structure}
% The table of blocks. Generate it from software/export/manifest.json rather
% than copying it from the paper - this thesis runs a 4-class variant, and the
% classifier is \thClasses{} outputs, not 1000.

\section{Quantization of Neural Networks}
\label{sec:quantization}

\subsection{Uniform Quantization}
\label{sec:uniform-quant}
% [jacob2018quantization], [krishnamoorthi2018whitepaper]. Define scale and
% zero-point once, here, in the notation \cref{sec:quant-scheme} will use.

\subsection{Per-Tensor and Per-Channel Quantization}
\label{sec:per-channel}

\subsection{Batch Normalization Folding}
\label{sec:bn-folding}
% [ioffe2015batchnorm] for BN itself; folding is what produces the large
% per-channel biases that broke ACC_W in \DD{019}. Set that up here and pay it
% off in \cref{sec:acc-bound}.

\subsection{Integer-Only Inference and Requantization}
\label{sec:requant-theory}
% The fixed-point multiplier (m0, shift) and rounding. This is the single most
% load-bearing subsection in the chapter: \cref{sec:requant-lanes} implements
% it and \cref{ch:verification} checks it bit for bit, so the rounding rule
% must be stated exactly, not approximately.

\section{Field-Programmable Gate Arrays}
\label{sec:fpga}

\subsection{FPGA Architecture}
\label{sec:fpga-arch}
% [betz1999architecture]. Brief.

\subsection{DSP48E2 Slices}
\label{sec:dsp48}
% [xilinx_ug579]. The internal pipeline registers matter to \cref{ch:timing}
% and the int8 packing question matters to \DD{021} - cover both, skip the rest.

\subsection{Block RAM and UltraRAM}
\label{sec:bram-uram}
% [xilinx_ug573]. URAM read latency is the reason 250 MHz was not reached
% (\DD{023}); give it enough space that \cref{sec:dd023} can just refer back.

\subsection{FPGA CAD Flow and Timing Analysis}
\label{sec:cad-flow}
% [xilinx_ug901]. Define WNS, slack and Fmax here, and define how this thesis
% computes Fmax from WNS: Fmax = 1000/(T - WNS). State the caveat that for a
% MET constraint this is a capability and the operating clock is the
% constrained one - \cref{sec:board-impl} depends on that distinction.

\subsection{Hardware Description Languages}
\label{sec:hdl}
% [ieee2001verilog]. Why Verilog-2001 specifically.

\subsection{High-Level Synthesis}
\label{sec:hls}
% Present fairly, then say why this thesis did not use it. \cref{sec:positioning}
% and \cref{ch:timing} carry the argument; do not repeat it here.

\subsection{Comparison of FPGAs with CPUs, GPUs and ASICs}
\label{sec:fpga-vs}
% [sze2017efficient]. Keep to the energy-per-inference framing that
% \cref{sec:power} will return to.

\section{The Zynq UltraScale+ MPSoC}
\label{sec:zynq}

\subsection{Processing System and Programmable Logic}
\label{sec:ps-pl}
% [xilinx_ug1085], [xilinx_ds891]. Include the PL clocking arrangement: the
% grid is coarse (1500 MHz over an integer divider gives 150 / 166.7 / 187.5 /
% 214.3 / 250 MHz only), and requesting 200 silently yields 187.5. That fact
% is a result in \cref{sec:board-impl}, so it needs a home in the background.

\subsection{The AXI Protocol Family}
\label{sec:axi}
% [arm_axi], [arm_axistream], [xilinx_pg021]. Cover exactly what wgt_dma and
% axi_lite_csr use: INCR bursts, the 4 KB boundary rule, the 256-beat cap,
% outstanding transactions, and AXI-Lite's error responses. Nothing else.

\section{Design Principles of CNN Accelerators}
\label{sec:accel-principles}

\subsection{Loop Tiling and Unrolling}
\label{sec:tiling}
% [zhang2015optimizing], [ma2017optimizing]. Define Tm, Tn, Tc in the notation
% \cref{sec:pw-array} uses, so DD-013 (32x16) and DD-014 (Tc=16) can be stated
% as choices in that space rather than re-derived.

\subsection{The Roofline Model}
\label{sec:roofline}
% [williams2009roofline]. docs/controller_design.md already contains the
% roofline for this design; \cref{sec:roofline-analysis} presents the result,
% so keep this to the model itself.
